\chapter{Boundaries as part of the calculus}
\label{ch:calculus}
\chapterintro{The constructions identify specific obligations for realization: compatibility, admissibility, generation, and effective evaluation. A failure at one level need not be a failure at another.}

\section{The ambient object and its admissible realizations}
The chain
\[
D\subseteq G_A\subseteq P
\]
contains three different mathematical roles. The direct sum provides the finite-support corrections. The product supplies every compatible prime-coordinate tuple. The intermediate group is the intended admissible realization, determined by a rational line in the quotient.

This arrangement gives a concrete model for a central question of Komorebi Space. Once the compatible ambient object exists, what additional data identifies the realization under study? Here the answer is exact: a rational subspace of $P/D$, or in the paired family its selector class. Existence of the ambient product does not eliminate that choice.

The construction does not assert that every Komorebi Space model is a prime-coordinate product. It isolates a theorem-level phenomenon that any broad realization language must be able to express: the space of compatible observations can be larger than the intended source image.

\section{Three certificates and one global obstruction}
The finite-seed argument supplies a projection certificate. An omitted prime coordinate vanishes on every admissible endomorphism orbit but not on the target group. This is a structural failure of generation under the fixed operator action.

The root algorithm supplies a coordinate certificate of impossibility or a finite correction witnessing success. It is a fully local algebraic realization problem. The successful and unsuccessful branches are both explicit.

The nonsplitting result supplies an extension obstruction. A rational quotient exists, but no homomorphic section realizes it inside the reduced product. In cocycle coordinates, the obstruction is the impossibility of removing the defect by a $D$-valued correction.

The grammar theorem is different. It excludes every finite description in a fixed language by an infinite diagonal comparison. It need not supply a terminating rejection procedure for individual candidate descriptions. Conflating this global semantic obstruction with the finite root certificates would overstate its algorithmic content.

\section{A realization profile rather than a single score}
One scalar cannot capture all the boundaries exhibited here. The following profile retains the distinctions:
\[
\bigl(\text{seed rank},\ \text{grammar complexity},\
\text{least presentation degree},\ \text{observation type}\bigr).
\]
For every $G_A$, the seed rank is $\aleph_0$ and the observation type is the same. The least presentation degree varies with $A$. Grammar complexity depends on which selectors the chosen semantic language can name.

The profile is not proposed as a complete classification invariant for arbitrary mathematical objects. It is a disciplined way to avoid substituting one question for another. A theory may have exact coordinates without a finite generator, a finite formula without computable equality, or an effectively presented inverse-limit completion without a computable copy of the intended dense realization.

\section{What a machine-facing contract should retain}
A reusable construction record should specify its input type, its permitted parameters, its output object, and the sense in which the output is determined. It should also state whether it returns a concrete representation, an abstract isomorphism type, or an existence proof. When completion is used, the intended image and the completed ambient object should be separate fields.

For the present family, a compact exact record is:
\begin{center}
\begin{tabularx}{\textwidth}{@{}p{0.22\textwidth}X@{}}
\toprule
Field&Content\\\midrule
Ambient system&$D\subseteq P$ and the quotient $q:P\to P/D$.\\
Parameter&The selector $A$, with its declared oracle or finite semantic description.\\
Admissible image&$q^{-1}(\Q q(u_A))$.\\
Concrete names&Reduced rationals paired with normalized finite-support vectors.\\
Operation&The finite cocycle law \eqref{eq:grouplaw}.\\
Positive certificate&A root and its exact multiplication check.\\
Negative certificate&A prime coordinate obstructing the required root.\\
Global invariants&Selector modulo finite changes and least presentation degree.\\
Observation boundary&Unmarked finite quotients forget the admissible source image.\\\bottomrule
\end{tabularx}
\end{center}
Such a record does not solve arbitrary group-theoretic problems. It makes the semantics of this construction inspectable and prevents a machine from promoting a finite observation theorem into a source reconstruction theorem.

\section{Further questions with exact starting points}
The general correspondence of Chapter~\ref{ch:purity} extends the ambient construction to higher-dimensional rational subspaces $W\le P/D$. Its finite quotient theorem is already complete. What remains is to classify those subspaces under the diagonal multipliers that induce group homomorphisms. The paired rank-one family has an exact answer, but the book does not claim that its binary signature classifies higher rank.

A second question concerns effective selection. For $G_A$, the normal-form operation and the paired root tests give a matching upper and lower bound. For a more general countable $W$, one must specify how its rational dependencies are presented and how lifts are computed. A list of candidate spanning vectors is not enough if their dependence relations are unavailable.

A third question concerns enriched observations. A coherent marked point restores the selector in the declared rank-one family. A useful general theorem would characterize which additional marks or admissibility predicates suffice for reconstruction in larger classes. Such a result must state whether its reconstruction is semantic, effective, or functorial.

Finally, the formal companion keeps the arithmetic kernel, quotient lemmas, tail signature, and computability reduction as distinct proof obligations. Appendix~\ref{app:formal} explains the resulting interfaces. The oracle theorem has its own effective semantics, with native oracle programs, computable equality, and terminating searches in addition to the algebraic proofs.

\section{Conclusion}
The principal boundary is now concrete. There is a fixed ambient completion, presented by a computable inverse system, and a family of countable, residually finite groups densely realized inside it. Their finite-quotient systems are canonically identified, and restriction to $D$ naturally identifies their homomorphism sets into every finite target. Yet the groups have distinct algebraic types and require arbitrarily different computational information.

Realization is therefore not exhausted by compatibility. It includes the selection of admissible objects, the transport of their meaning, and an account of the information needed to evaluate them. The groups in this book turn that distinction into explicit constructions, exact classifications, and reproducible finite certificates.
